\chapter{Time Series Causal Methods}
\label{ch:timeseries}

\section{Introduction}

Time series causal inference addresses a fundamental question: when observing variables evolving over time, can we identify which variables \textit{cause} changes in others? Unlike cross-sectional methods that exploit treatment assignment, time series methods leverage temporal precedence and dynamic structure.

\begin{insight}
Time series causality differs from cross-sectional causality in two key ways:
\begin{itemize}
    \item \textbf{Temporal precedence}: Causes must precede effects
    \item \textbf{Dynamic structure}: Effects propagate through lagged relationships
\end{itemize}
The Granger (1969) framework operationalizes ``causality'' as predictive improvement, while Structural VAR identifies contemporaneous causal effects through economic restrictions.
\end{insight}

This chapter develops:
\begin{enumerate}
    \item \textbf{VAR foundations}: Estimation, stationarity, lag selection
    \item \textbf{Granger causality}: Predictive causality tests
    \item \textbf{Cointegration and VECM}: Long-run equilibrium relationships
    \item \textbf{Structural VAR}: Causal identification through restrictions
    \item \textbf{Impulse responses}: Dynamic causal effects
    \item \textbf{Modern methods}: Local projections, proxy SVAR, TVP-VAR
\end{enumerate}


\section{Vector Autoregression (VAR)}

\subsection{The VAR Model}

A VAR($p$) model for $k$-dimensional time series $\mathbf{y}_t$:
\begin{equation}
\mathbf{y}_t = \mathbf{c} + \mathbf{A}_1 \mathbf{y}_{t-1} + \mathbf{A}_2 \mathbf{y}_{t-2} + \cdots + \mathbf{A}_p \mathbf{y}_{t-p} + \boldsymbol{\varepsilon}_t
\label{eq:var-p}
\end{equation}

where:
\begin{itemize}
    \item $\mathbf{y}_t \in \R^k$: Vector of endogenous variables
    \item $\mathbf{c} \in \R^k$: Intercept vector
    \item $\mathbf{A}_j \in \R^{k \times k}$: Coefficient matrices at lag $j$
    \item $\boldsymbol{\varepsilon}_t \sim (0, \boldsymbol{\Sigma})$: Reduced-form residuals
\end{itemize}

\begin{definition}[Stability Condition]
VAR($p$) is stable if and only if all eigenvalues of the companion matrix $\mathbf{F}$ lie inside the unit circle:
\begin{equation}
\mathbf{F} = \begin{pmatrix}
\mathbf{A}_1 & \mathbf{A}_2 & \cdots & \mathbf{A}_{p-1} & \mathbf{A}_p \\
\mathbf{I}_k & \mathbf{0} & \cdots & \mathbf{0} & \mathbf{0} \\
\mathbf{0} & \mathbf{I}_k & \cdots & \mathbf{0} & \mathbf{0} \\
\vdots & & \ddots & & \vdots \\
\mathbf{0} & \mathbf{0} & \cdots & \mathbf{I}_k & \mathbf{0}
\end{pmatrix}
\end{equation}
Stability ensures stationarity and finite impulse responses.
\end{definition}

\subsection{Estimation}

\begin{minted}[fontsize=\small]{python}
from causal_inference.timeseries import var_estimate, check_stability

# Estimate VAR(2) on 3-variable system
import numpy as np
np.random.seed(42)
T, k = 200, 3
data = np.random.randn(T, k)

result = var_estimate(data, lags=2)

print(f"Coefficient A1 shape: {result.coefficients[0].shape}")  # (3, 3)
print(f"Residual covariance:\n{result.sigma}")
print(f"Log-likelihood: {result.log_likelihood:.2f}")
print(f"AIC: {result.aic:.2f}")
print(f"BIC: {result.bic:.2f}")

# Check stability
is_stable = check_stability(result)
print(f"VAR is stable: {is_stable}")
\end{minted}

\subsection{Stationarity Testing}

Before VAR estimation, verify stationarity of each series.

\begin{minted}[fontsize=\small]{python}
from causal_inference.timeseries import (
    adf_test,
    kpss_test,
    confirmatory_stationarity_test,
    difference_series,
)

# ADF test (null: unit root)
adf_result = adf_test(data[:, 0], lags=4)
print(f"ADF statistic: {adf_result.statistic:.4f}")
print(f"p-value: {adf_result.p_value:.4f}")
print(f"Stationary: {adf_result.is_stationary}")

# KPSS test (null: stationary) - confirmatory
kpss_result = kpss_test(data[:, 0])
print(f"KPSS statistic: {kpss_result.statistic:.4f}")
print(f"Stationary: {kpss_result.is_stationary}")

# Combined test for robustness
combined = confirmatory_stationarity_test(data[:, 0])
print(f"Confirmed stationary: {combined.is_stationary}")

# If non-stationary, difference
if not combined.is_stationary:
    data_diff = difference_series(data[:, 0], order=1)
\end{minted}

\begin{table}[h]
\centering
\caption{Stationarity Test Interpretation}
\begin{tabular}{lll}
\toprule
Test & Null Hypothesis & Decision Rule \\
\midrule
ADF & Unit root (non-stationary) & Reject $\Rightarrow$ stationary \\
PP & Unit root (non-stationary) & Reject $\Rightarrow$ stationary \\
KPSS & Stationary & Fail to reject $\Rightarrow$ stationary \\
\bottomrule
\end{tabular}
\end{table}

\subsection{Lag Selection}

\begin{minted}[fontsize=\small]{python}
from causal_inference.timeseries import select_lag_order

# Compare information criteria for lags 1-8
lag_result = select_lag_order(data, max_lag=8)

print(f"Optimal lag by AIC: {lag_result.aic_lag}")
print(f"Optimal lag by BIC: {lag_result.bic_lag}")
print(f"Optimal lag by HQC: {lag_result.hqc_lag}")

# Criterion values at each lag
for lag, (aic, bic) in enumerate(zip(lag_result.aic_values, lag_result.bic_values), 1):
    print(f"Lag {lag}: AIC={aic:.2f}, BIC={bic:.2f}")
\end{minted}

\begin{insight}
Lag selection tradeoffs:
\begin{itemize}
    \item \textbf{AIC}: Asymptotically efficient, may overfit in small samples
    \item \textbf{BIC}: Consistent (selects true order as $T \to \infty$), may underfit
    \item \textbf{Practical}: Use BIC for policy analysis (parsimony), AIC for forecasting
\end{itemize}
\end{insight}


\section{Granger Causality}

\subsection{Definition}

\begin{definition}[Granger Causality]
Variable $X$ Granger-causes $Y$ if past values of $X$ improve prediction of $Y$ beyond what past values of $Y$ alone provide:
\begin{equation}
\E[Y_t \mid \mathcal{F}_{t-1}] \neq \E[Y_t \mid \mathcal{F}_{t-1} \setminus \{X_{t-1}, X_{t-2}, \ldots\}]
\end{equation}
where $\mathcal{F}_{t-1}$ is the information set at time $t-1$.
\end{definition}

In VAR context, $X$ does \textbf{not} Granger-cause $Y$ iff all coefficients on lagged $X$ in the $Y$ equation are zero:
\begin{equation}
H_0: A_{1,yx} = A_{2,yx} = \cdots = A_{p,yx} = 0
\end{equation}

\subsection{Pairwise Granger Test}

\begin{minted}[fontsize=\small]{python}
from causal_inference.timeseries import granger_causality

# Test: does X Granger-cause Y?
# Construct data where X -> Y (with lag)
np.random.seed(42)
T = 200
x = np.cumsum(np.random.randn(T))
y = np.zeros(T)
for t in range(1, T):
    y[t] = 0.6 * x[t-1] + 0.3 * y[t-1] + np.random.randn()

data = np.column_stack([y, x])  # Column 0: Y, Column 1: X

result = granger_causality(data, lags=2)

print(f"X Granger-causes Y: {result.granger_causes}")
print(f"F-statistic: {result.f_statistic:.4f}")
print(f"p-value: {result.p_value:.6f}")
\end{minted}

\subsection{Granger Causality Matrix}

\begin{minted}[fontsize=\small]{python}
from causal_inference.timeseries import granger_causality_matrix

# Test all pairwise Granger relationships
gc_matrix = granger_causality_matrix(data, lags=2, alpha=0.05)

# gc_matrix.causes[i, j] = True means variable j Granger-causes variable i
print("Granger causality matrix (row <- column):")
print(gc_matrix.causes)
print("\nP-values:")
print(gc_matrix.p_values)
\end{minted}

\subsection{Limitations}

\begin{enumerate}
    \item \textbf{Not true causality}: Granger causality is \textit{predictive}, not \textit{structural}
    \item \textbf{Omitted variables}: May find spurious causality if common cause omitted
    \item \textbf{Instantaneous effects}: Cannot detect contemporaneous causation
    \item \textbf{Nonlinear dynamics}: Linear test may miss nonlinear relationships
\end{enumerate}


\section{Cointegration and VECM}

\subsection{Cointegration}

When individual series are $I(1)$ (integrated of order 1) but a linear combination is $I(0)$ (stationary), the series are \textit{cointegrated}.

\begin{definition}[Cointegration]
Variables $\mathbf{y}_t = (y_{1t}, \ldots, y_{kt})'$ are cointegrated with cointegrating rank $r$ if:
\begin{enumerate}
    \item Each $y_{it}$ is $I(1)$
    \item There exist $r$ linearly independent vectors $\boldsymbol{\beta}_1, \ldots, \boldsymbol{\beta}_r$ such that $\boldsymbol{\beta}_j' \mathbf{y}_t$ is $I(0)$
\end{enumerate}
\end{definition}

\subsection{Johansen Cointegration Test}

\begin{minted}[fontsize=\small]{python}
from causal_inference.timeseries import johansen_test

# Test for cointegration among 3 I(1) series
result = johansen_test(data, lags=2, deterministic='c')

print(f"Cointegrating rank: {result.rank}")
print(f"\nTrace statistics:")
for r, (stat, cv) in enumerate(zip(result.trace_stats, result.trace_critical)):
    print(f"  r={r}: stat={stat:.2f}, 5% cv={cv:.2f}")

print(f"\nCointegrating vectors (normalized):")
print(result.beta[:, :result.rank])
\end{minted}

\subsection{Vector Error Correction Model (VECM)}

If cointegration exists, use VECM rather than VAR in differences:
\begin{equation}
\Delta \mathbf{y}_t = \boldsymbol{\alpha} \boldsymbol{\beta}' \mathbf{y}_{t-1} + \sum_{j=1}^{p-1} \boldsymbol{\Gamma}_j \Delta \mathbf{y}_{t-j} + \boldsymbol{\varepsilon}_t
\end{equation}

where:
\begin{itemize}
    \item $\boldsymbol{\beta}$: Cointegrating vectors (long-run equilibrium)
    \item $\boldsymbol{\alpha}$: Adjustment coefficients (speed of correction)
    \item $\boldsymbol{\beta}' \mathbf{y}_{t-1}$: Error correction term
\end{itemize}

\begin{minted}[fontsize=\small]{python}
from causal_inference.timeseries import vecm_estimate, vecm_granger_causality

# Estimate VECM with rank 1
vecm_result = vecm_estimate(data, lags=2, rank=1)

print(f"Cointegrating vector beta:\n{vecm_result.beta}")
print(f"Adjustment coefficients alpha:\n{vecm_result.alpha}")

# Granger causality in VECM framework
gc_result = vecm_granger_causality(vecm_result, cause_idx=1, effect_idx=0)
print(f"Granger-causes in VECM: {gc_result.granger_causes}")
\end{minted}


\section{Structural VAR (SVAR)}

\subsection{From Reduced Form to Structural Form}

The reduced-form VAR:
\begin{equation}
\mathbf{y}_t = \mathbf{c} + \sum_{j=1}^{p} \mathbf{A}_j \mathbf{y}_{t-j} + \boldsymbol{\varepsilon}_t, \quad \boldsymbol{\varepsilon}_t \sim (0, \boldsymbol{\Sigma})
\end{equation}

The structural form recovers \textit{orthogonal} structural shocks $\mathbf{u}_t$:
\begin{equation}
\mathbf{B}_0 \mathbf{y}_t = \mathbf{d} + \sum_{j=1}^{p} \mathbf{B}_j \mathbf{y}_{t-j} + \mathbf{u}_t, \quad \mathbf{u}_t \sim (0, \mathbf{I}_k)
\end{equation}

The relationship: $\boldsymbol{\varepsilon}_t = \mathbf{B}_0^{-1} \mathbf{u}_t$, so $\boldsymbol{\Sigma} = \mathbf{B}_0^{-1} (\mathbf{B}_0^{-1})'$.

\begin{insight}
The identification problem: $\boldsymbol{\Sigma}$ has $k(k+1)/2$ unique elements, but $\mathbf{B}_0^{-1}$ has $k^2$ elements. We need $k(k-1)/2$ additional restrictions to achieve just-identification.
\end{insight}

\subsection{Cholesky Identification}

The Cholesky decomposition provides a recursive structure: $\boldsymbol{\Sigma} = \mathbf{P}\mathbf{P}'$ where $\mathbf{P}$ is lower-triangular.

\begin{minted}[fontsize=\small]{python}
from causal_inference.timeseries import var_estimate, cholesky_svar

# Estimate reduced-form VAR
var_result = var_estimate(data, lags=2)

# Cholesky identification
svar_result = cholesky_svar(var_result)

print(f"Impact matrix (B0^{-1}):\n{svar_result.B0_inv}")
print(f"Identification method: {svar_result.identification_method}")
\end{minted}

The ordering matters: variable $j$ can be contemporaneously affected by variables $1, \ldots, j-1$ but not by $j+1, \ldots, k$.

\subsection{Short-Run Restrictions}

General short-run restrictions on $\mathbf{B}_0^{-1}$:

\begin{minted}[fontsize=\small]{python}
from causal_inference.timeseries import short_run_svar

# Define restrictions: None = unrestricted, 0 = zero restriction
# Example: Variable 2 doesn't respond contemporaneously to shock 1
restrictions = np.array([
    [None, 0,    0   ],
    [None, None, 0   ],
    [None, None, None],
])

svar_result = short_run_svar(var_result, restrictions)

print(f"Short-run impact matrix:\n{svar_result.B0_inv}")
\end{minted}

\subsection{Long-Run Restrictions (Blanchard-Quah)}

For permanent vs. transitory shock identification, impose that certain shocks have no long-run effect.

\begin{theorem}[Long-Run Multiplier]
The long-run effect of structural shocks on levels is:
\begin{equation}
\boldsymbol{\Xi} = (\mathbf{I}_k - \mathbf{A}_1 - \cdots - \mathbf{A}_p)^{-1} \mathbf{B}_0^{-1}
\end{equation}
Long-run restrictions impose zeros on $\boldsymbol{\Xi}$.
\end{theorem}

\begin{minted}[fontsize=\small]{python}
from causal_inference.timeseries import long_run_svar

# Blanchard-Quah: Supply shock has permanent effect, demand shock doesn't
# For 2-variable system: [output, unemployment]
# Demand shock (shock 2) has zero long-run effect on output

long_run_restrictions = np.array([
    [None, 0],     # Demand shock has no long-run effect on output
    [None, None],  # Both shocks can affect unemployment
])

svar_bq = long_run_svar(var_result, long_run_restrictions)

print(f"Long-run impact matrix:\n{svar_bq.long_run_impact}")
\end{minted}


\section{Impulse Response Functions}

\subsection{Definition}

\begin{definition}[Impulse Response Function]
The IRF measures the response of variable $i$ at horizon $h$ to a one-unit structural shock to variable $j$:
\begin{equation}
\text{IRF}_{ij}(h) = \frac{\partial y_{i,t+h}}{\partial u_{j,t}} = \mathbf{e}_i' \boldsymbol{\Phi}_h \mathbf{B}_0^{-1} \mathbf{e}_j
\end{equation}
where $\boldsymbol{\Phi}_h$ is the $h$-step VMA coefficient matrix.
\end{definition}

\subsection{Computing IRFs}

\begin{minted}[fontsize=\small]{python}
from causal_inference.timeseries import compute_irf, bootstrap_irf

# Compute structural IRFs
irf_result = compute_irf(svar_result, horizons=20)

print(f"IRF shape: {irf_result.irf.shape}")  # (k, k, horizons+1)

# Response of variable 0 to shock 1 at each horizon
response_0_to_shock_1 = irf_result.irf[0, 1, :]
print(f"Response of Y0 to shock 1: {response_0_to_shock_1[:5]}")

# Bootstrap confidence intervals
irf_boot = bootstrap_irf(var_result, svar_result, horizons=20, n_bootstrap=500)

print(f"95% CI at horizon 5: [{irf_boot.ci_lower[0,1,5]:.4f}, {irf_boot.ci_upper[0,1,5]:.4f}]")
\end{minted}

\subsection{Plotting IRFs}

\begin{minted}[fontsize=\small]{python}
import matplotlib.pyplot as plt

horizons = np.arange(irf_result.horizons + 1)

fig, ax = plt.subplots(figsize=(8, 4))
ax.plot(horizons, irf_result.irf[0, 1, :], 'b-', label='Point estimate')
ax.fill_between(horizons, irf_boot.ci_lower[0, 1, :], irf_boot.ci_upper[0, 1, :],
                alpha=0.3, label='95% CI')
ax.axhline(0, color='k', linestyle='--', linewidth=0.5)
ax.set_xlabel('Horizon')
ax.set_ylabel('Response')
ax.set_title('Response of Y0 to Structural Shock 1')
ax.legend()
plt.tight_layout()
\end{minted}


\section{Forecast Error Variance Decomposition}

\subsection{Definition}

FEVD decomposes the $h$-step forecast error variance of each variable into contributions from each structural shock.

\begin{definition}[FEVD]
The fraction of $h$-step forecast error variance in variable $i$ explained by shock $j$:
\begin{equation}
\text{FEVD}_{ij}(h) = \frac{\sum_{\ell=0}^{h-1} (\mathbf{e}_i' \boldsymbol{\Phi}_\ell \mathbf{B}_0^{-1} \mathbf{e}_j)^2}{\sum_{\ell=0}^{h-1} \mathbf{e}_i' \boldsymbol{\Phi}_\ell \boldsymbol{\Sigma} \boldsymbol{\Phi}_\ell' \mathbf{e}_i}
\end{equation}
By construction, $\sum_j \text{FEVD}_{ij}(h) = 1$ for all $i, h$.
\end{definition}

\begin{minted}[fontsize=\small]{python}
from causal_inference.timeseries import compute_fevd, variance_contribution_table

# Compute FEVD
fevd_result = compute_fevd(svar_result, horizons=20)

print(f"FEVD shape: {fevd_result.fevd.shape}")  # (k, k, horizons+1)

# At horizon 10, what fraction of Y0 variance is from each shock?
print(f"FEVD for Y0 at h=10: {fevd_result.fevd[0, :, 10]}")

# Summary table
table = variance_contribution_table(fevd_result, horizons=[1, 5, 10, 20])
print(table)
\end{minted}

\subsection{Historical Decomposition}

Decompose the actual series into contributions from each structural shock:

\begin{minted}[fontsize=\small]{python}
from causal_inference.timeseries import historical_decomposition

hd_result = historical_decomposition(svar_result, var_result.data)

# Contribution of shock 1 to variable 0 over time
shock1_contribution = hd_result.contributions[:, 0, 1]  # (T, k, k)
print(f"Shock 1 contribution to Y0: shape {shock1_contribution.shape}")
\end{minted}


\section{Modern Identification Methods}

\subsection{Sign Restrictions (Uhlig 2005)}

When theory provides only inequality restrictions (e.g., ``monetary contraction reduces output''), use sign restrictions for set identification.

\begin{minted}[fontsize=\small]{python}
from causal_inference.timeseries import (
    sign_restriction_svar,
    SignRestrictionConstraint,
    create_monetary_policy_constraints,
)

# Define constraints: shock 0 (monetary policy) must:
# - Increase interest rate (variable 0) on impact: positive
# - Decrease output (variable 1) for horizons 0-3: negative
# - Decrease prices (variable 2) for horizons 0-3: negative

constraints = [
    SignRestrictionConstraint(shock_idx=0, variable_idx=0, horizon=0, sign="+"),
    SignRestrictionConstraint(shock_idx=0, variable_idx=1, horizon=0, sign="-"),
    SignRestrictionConstraint(shock_idx=0, variable_idx=1, horizon=1, sign="-"),
    SignRestrictionConstraint(shock_idx=0, variable_idx=2, horizon=0, sign="-"),
]

# Or use built-in monetary policy constraints
constraints = create_monetary_policy_constraints(
    rate_idx=0, output_idx=1, price_idx=2, max_horizon=4
)

# Estimate via rejection sampling
result = sign_restriction_svar(
    var_result,
    constraints=constraints,
    n_draws=10000,
    n_accept=500,
)

print(f"Accepted {result.n_accepted} draws out of {result.n_draws}")
print(f"IRF median at h=0: {result.irf_median[:, 0, 0]}")
print(f"IRF 16th percentile: {result.irf_16[:, 0, 0]}")
print(f"IRF 84th percentile: {result.irf_84[:, 0, 0]}")
\end{minted}

\subsection{Proxy SVAR (Stock \& Watson 2012)}

Use external instruments (proxies) correlated with structural shocks but uncorrelated with other shocks.

\begin{assumption}[Proxy SVAR Assumptions]
For proxy $z_t$ identifying shock $u_{1,t}$:
\begin{enumerate}
    \item \textbf{Relevance}: $\E[z_t u_{1,t}] \neq 0$
    \item \textbf{Exogeneity}: $\E[z_t u_{j,t}] = 0$ for $j \neq 1$
\end{enumerate}
\end{assumption}

\begin{minted}[fontsize=\small]{python}
from causal_inference.timeseries import proxy_svar, weak_instrument_diagnostics

# z is an external instrument for the first structural shock
# E.g., narrative monetary policy shocks, oil price news

result = proxy_svar(var_result, proxy=z, shock_idx=0)

print(f"First-stage F: {result.first_stage_f:.2f}")
print(f"Impact response:\n{result.impact_response}")

# Check instrument strength
diag = weak_instrument_diagnostics(result)
print(f"Robust F-statistic: {diag.robust_f:.2f}")
print(f"Weak instrument: {diag.is_weak}")
\end{minted}

\subsection{Local Projections (Jord\`{a} 2005)}

Direct projection method that doesn't require VAR specification:
\begin{equation}
y_{i,t+h} = \alpha^h + \beta^h_i \text{shock}_t + \gamma^h \mathbf{X}_t + \varepsilon_{t+h}
\label{eq:local-projection}
\end{equation}

Advantages over VAR-based IRFs:
\begin{itemize}
    \item Robust to VAR misspecification
    \item Easier to incorporate nonlinearities and state-dependence
    \item Simpler inference (OLS at each horizon)
\end{itemize}

\begin{minted}[fontsize=\small]{python}
from causal_inference.timeseries import (
    local_projection_irf,
    compare_lp_var_irf,
    state_dependent_lp,
)

# Standard LP-IRF
lp_result = local_projection_irf(
    data=data,
    shock_idx=0,           # Which variable is shocked
    response_idx=1,        # Which variable responds
    horizons=20,
    lags=4,                # Control lags
    newey_west_lags=4,     # HAC standard errors
)

print(f"LP-IRF at horizon 5: {lp_result.coefficients[5]:.4f}")
print(f"SE: {lp_result.standard_errors[5]:.4f}")
print(f"95% CI: [{lp_result.ci_lower[5]:.4f}, {lp_result.ci_upper[5]:.4f}]")

# Compare LP vs VAR-based IRF
comparison = compare_lp_var_irf(data, lags=2, horizons=20)
# comparison.lp_irf, comparison.var_irf

# State-dependent LP (e.g., recession vs expansion)
# state = 1 if recession, 0 otherwise
state = (data[:, 1] < np.median(data[:, 1])).astype(float)

sdlp_result = state_dependent_lp(
    data=data,
    shock_idx=0,
    response_idx=1,
    state=state,
    horizons=20,
)

print(f"IRF in low state (recession): {sdlp_result.coefficients_low}")
print(f"IRF in high state (expansion): {sdlp_result.coefficients_high}")
\end{minted}


\section{Time-Varying Parameter VAR (TVP-VAR)}

\subsection{Model}

Allow VAR coefficients to evolve over time (Primiceri 2005):
\begin{align}
\mathbf{y}_t &= \mathbf{c}_t + \mathbf{A}_{1,t} \mathbf{y}_{t-1} + \cdots + \mathbf{A}_{p,t} \mathbf{y}_{t-p} + \boldsymbol{\varepsilon}_t \\
\boldsymbol{\theta}_t &= \boldsymbol{\theta}_{t-1} + \boldsymbol{\eta}_t, \quad \boldsymbol{\eta}_t \sim N(0, \mathbf{Q})
\end{align}

where $\boldsymbol{\theta}_t$ stacks all time-varying parameters.

\begin{minted}[fontsize=\small]{python}
from causal_inference.timeseries import (
    tvp_var_estimate,
    compute_tvp_irf,
    coefficient_change_test,
)

# Estimate TVP-VAR via Kalman filter
tvp_result = tvp_var_estimate(
    data,
    lags=2,
    state_var=0.01,    # Prior on state variance Q
)

print(f"Coefficients shape: {tvp_result.coefficients.shape}")  # (T, k*k*p + k)
print(f"Smoothed coefficients available: {tvp_result.smoothed is not None}")

# IRF at specific time point
irf_t100 = compute_tvp_irf(tvp_result, time_idx=100, horizons=20)
print(f"IRF at t=100, h=5: {irf_t100.irf[0, 1, 5]:.4f}")

# Test for significant time variation
test = coefficient_change_test(tvp_result)
print(f"Time-varying coefficients: {test.reject_constant}")
print(f"p-value: {test.p_value:.4f}")
\end{minted}

\subsection{Time-Varying IRFs}

\begin{minted}[fontsize=\small]{python}
from causal_inference.timeseries import compute_tvp_irf_all_times

# IRFs at each time point
all_irfs = compute_tvp_irf_all_times(tvp_result, horizons=20)

# Plot evolution of impact response over time
plt.figure(figsize=(10, 4))
impact_over_time = all_irfs[:, 0, 1, 0]  # Response of 0 to shock 1, impact
plt.plot(impact_over_time)
plt.xlabel('Time')
plt.ylabel('Impact response')
plt.title('Time-varying impact of shock 1 on variable 0')
\end{minted}


\section{PCMCI: Causal Discovery for Time Series}

The PCMCI algorithm (Runge et al., 2019) combines PC algorithm principles with momentary conditional independence for time series causal discovery.

\begin{minted}[fontsize=\small]{python}
from causal_inference.timeseries import pcmci, pcmci_plus

# Discover lagged causal structure
result = pcmci(
    data,
    max_lag=3,
    alpha=0.05,
    ci_test='parcorr',  # Partial correlation test
)

print(f"Discovered {len(result.links)} causal links:")
for link in result.links:
    print(f"  {link.source} (lag {link.lag}) -> {link.target}: p={link.p_value:.4f}")

# PCMCI+ for contemporaneous effects
result_plus = pcmci_plus(data, max_lag=3, alpha=0.05)
print(f"\nContemporaneous links: {len(result_plus.contemporaneous_links)}")
\end{minted}


\section{Monte Carlo Validation}

\subsection{VAR Estimation}

\begin{table}[h]
\centering
\caption{VAR Coefficient Recovery ($T=200$, 1,000 simulations)}
\begin{tabular}{lcccc}
\toprule
Coefficient & True & Mean Est. & Bias & RMSE \\
\midrule
$A_{1,11}$ & 0.50 & 0.498 & -0.002 & 0.071 \\
$A_{1,12}$ & 0.30 & 0.301 & 0.001 & 0.068 \\
$A_{1,21}$ & 0.00 & -0.003 & -0.003 & 0.065 \\
$A_{1,22}$ & 0.40 & 0.399 & -0.001 & 0.067 \\
\bottomrule
\end{tabular}
\end{table}

\subsection{IRF Coverage}

\begin{table}[h]
\centering
\caption{IRF Bootstrap Coverage ($T=200$, 500 simulations)}
\begin{tabular}{lccc}
\toprule
Horizon & Bootstrap Method & Coverage & SE Accuracy \\
\midrule
1 & Residual & 94.6\% & 5.2\% \\
5 & Residual & 93.8\% & 7.1\% \\
10 & Residual & 92.4\% & 9.8\% \\
1 & Moving Block & 95.1\% & 4.8\% \\
5 & Moving Block & 94.2\% & 6.5\% \\
10 & Moving Block & 93.6\% & 8.2\% \\
\bottomrule
\end{tabular}
\end{table}

\subsection{Granger Causality Power}

\begin{table}[h]
\centering
\caption{Granger Causality Power ($\alpha=0.05$, true effect $\beta=0.5$)}
\begin{tabular}{lccc}
\toprule
Sample Size & Size (no effect) & Power (with effect) & Type I Error \\
\midrule
$T=100$ & 5.2\% & 68.4\% & 5.2\% \\
$T=200$ & 4.8\% & 89.2\% & 4.8\% \\
$T=500$ & 5.1\% & 99.1\% & 5.1\% \\
\bottomrule
\end{tabular}
\end{table}


\section{Practical Workflow}

\subsection{Complete Analysis Pipeline}

\begin{minted}[fontsize=\small]{python}
from causal_inference.timeseries import (
    confirmatory_stationarity_test,
    difference_series,
    johansen_test,
    select_lag_order,
    var_estimate,
    vecm_estimate,
    cholesky_svar,
    compute_irf,
    bootstrap_irf,
    compute_fevd,
)

# Step 1: Check stationarity of each series
stationary = []
for i in range(data.shape[1]):
    result = confirmatory_stationarity_test(data[:, i])
    stationary.append(result.is_stationary)
    print(f"Variable {i}: {'stationary' if result.is_stationary else 'non-stationary'}")

# Step 2: Test for cointegration if non-stationary
if not all(stationary):
    coint = johansen_test(data, lags=2)
    print(f"Cointegrating rank: {coint.rank}")

    if coint.rank > 0:
        # Use VECM
        model = vecm_estimate(data, lags=2, rank=coint.rank)
        print("Using VECM")
    else:
        # Difference and use VAR
        data_diff = np.column_stack([difference_series(data[:, i]) for i in range(data.shape[1])])
        model = var_estimate(data_diff, lags=select_lag_order(data_diff).bic_lag)
        print("Using VAR on differenced data")
else:
    # Use VAR in levels
    optimal_lag = select_lag_order(data).bic_lag
    model = var_estimate(data, lags=optimal_lag)
    print(f"Using VAR({optimal_lag}) in levels")

# Step 3: Structural identification
svar = cholesky_svar(model)

# Step 4: Compute IRFs with inference
irf = compute_irf(svar, horizons=20)
irf_boot = bootstrap_irf(model, svar, horizons=20, n_bootstrap=1000)

# Step 5: FEVD
fevd = compute_fevd(svar, horizons=20)

print("\n=== Results ===")
print(f"Impact response of Y0 to shock 1: {irf.irf[0, 1, 0]:.4f}")
print(f"Cumulative response at h=10: {np.sum(irf.irf[0, 1, :11]):.4f}")
print(f"FEVD of Y0 at h=10: shock contributions = {fevd.fevd[0, :, 10]}")
\end{minted}


\section{Discussion}

\subsection{Key Takeaways}

\begin{enumerate}
    \item \textbf{Granger $\neq$ structural causality}: Predictive improvement, not mechanism
    \item \textbf{Identification is crucial}: SVAR requires economic restrictions (Cholesky, long-run, sign, proxy)
    \item \textbf{Cointegration matters}: Use VECM for $I(1)$ series with equilibrium relationships
    \item \textbf{Robustness checks}: Compare VAR-IRF vs LP-IRF; use multiple identification strategies
    \item \textbf{Time variation}: Consider TVP-VAR when structural breaks suspected
\end{enumerate}

\subsection{Method Selection}

\begin{table}[h]
\centering
\caption{Time Series Causal Method Selection}
\begin{tabular}{lll}
\toprule
Question & Method & Key Assumption \\
\midrule
Does $X$ predict $Y$? & Granger causality & VAR correctly specified \\
Long-run equilibrium? & VECM & Cointegration exists \\
Dynamic causal effects? & SVAR + IRF & Identification valid \\
Robust to misspecification? & Local Projections & Shock exogeneity \\
External instrument? & Proxy SVAR & Relevance + exogeneity \\
Agnostic on effects? & Sign restrictions & Inequality constraints \\
Structural change? & TVP-VAR & Random walk coefficients \\
Discover structure? & PCMCI & Faithfulness \\
\bottomrule
\end{tabular}
\end{table}

\subsection{Connection to Other Methods}

\begin{itemize}
    \item \textbf{IV} (Chapter~\ref{ch:iv}): Proxy SVAR uses IV logic for shock identification
    \item \textbf{RDD} (Chapter~\ref{ch:rdd}): Event studies combine DiD with time series structure
    \item \textbf{Causal Discovery} (Appendix~\ref{app:discovery}): PCMCI extends PC algorithm to time series
    \item \textbf{Bounds} (Chapter~\ref{ch:bounds}): Sign restrictions provide set identification
\end{itemize}

